% -------------------------------------------------------------------------
% ÜBERHOLT (31.08.2026). Maßgeblich ist die eingereichte main.tex. Zahlen,
% Zeilennummern und Angaben zum Overleaf-Stand in dieser Datei stammen aus
% der Zeit VOR der Crosswalk-Korrektur vom 29.08.2026 (35 statt 36 Stadt-
% teile, vollständiger Neulauf) und wurden bewusst NICHT nachgezogen: die
% Datei ist als datierter Entwurf Teil der Arbeitsdokumentation.
% -------------------------------------------------------------------------
% =========================================================================
% KAPITEL 5.4 -- ANALYSERAHMEN UND BASELINES
% Fassung 1, 25.08.2026. Zweiter Schwerpunkt des Kapitels.
%
% ERSETZT den Platzhalterabsatz in main.tex Z. 1582-1584 ("Der zweite
% Schwerpunkt. Er legt fest, wie die Datensaetze in Trainings- und
% Testmengen zerfallen ...").
%
% GLEITOBJEKTE: docs/kapitel5-4_gleitobjekte.tex
%   Abbildung  abb:foldstruktur   -> nach Absatz 2
%   Quelltext  lst:foldmasken     -> nach Absatz 2      [KUERZBAR]
%   Quelltext  lst:folds          -> nach Absatz 4
%   Quelltext  lst:merkmalsliste  -> nach Absatz 8      [KUERZBAR]
%   Quelltext  lst:dtypes         -> nach Absatz 9
%   Tabelle    tab:baselines-menge     -> nach Absatz 11
%   Quelltext  lst:poisson        -> nach Absatz 11
%   Tabelle    tab:baselines-struktur  -> nach Absatz 12
%   Quelltext  lst:logit          -> nach Absatz 12
%   Quelltext  lst:test-folds     -> nach Absatz 16     [KUERZBAR]
%   Tabelle    tab:steckbrief     -> ans ENDE des Abschnitts
%
% NEUE ZITATE: Chapman2000 (A1), Roberts2017 (A2), Killip2004 (A7),
% Gourieroux1981 (A13). Chapman, Roberts und Gourieroux stehen bereits in
% literatur.bib; Killip2004 muss noch eingefuegt werden
% (docs/literatur_kapitel5-4_killip.bib). Roberts2017 und Gourieroux1981
% werden hier ZUM ERSTEN MAL zitiert -- vorher erschienen sie nicht im
% Literaturverzeichnis, weil biblatex nur zitierte Eintraege druckt.
%
% ERLEDIGTE AUFLAGEN
%   A (Abweichung vom Expose)  -> Absatz 3, Ausfuehrung in Kapitel 8
%   B ("Generalisierung auf unbekannte Stadtteile" woertlich) -> Absatz 1
%   C (identische Merkmale und Splits BELEGEN) -> Absatz 8
%   D (drei Argumente woertlich, in Schroeters Reihenfolge)   -> Absatz 14
%   R-2  (Vorbehalt zu Argument 3)      -> Absatz 14
%   R-14 (Wechsel weg von NegBin)       -> Absatz 13
%   R-5  (std_wiederholungen statt std_folds) -> Absatz 6
%   Serialisierung (Auflage 10.08.)     -> Absatz 9
%
% WAS BEWUSST NICHT DRINSTEHT
%   - class_weight="balanced" als ERKLAERUNG fuer das Verhaeltnis von
%     Trefferquote und Macro-F1 -> Kapitel 6. Hier steht nur der Befund
%     und die Konstruktionsentscheidung (keine Vervielfachung von Zeilen).
%   - Der flache Entscheidungsbaum (Macro-F1 0,270) -> Ergebnis, Kapitel 7.
%   - Der Extrapolationsanteil von 33,7 % -> Kapitel 7 und 8.
%
% BAUSTELLE, DIE HIER AUFFAELLT
%   In Kapitel 2.4 (sec:referenzmodelle) steht bislang KEIN Text, nur
%   Kommentare. Die Verweise unten zeigen deshalb auf 2.6 und 3.2, nicht
%   auf 2.4. Sobald 2.4 geschrieben ist, koennen Absatz 11 und 12 je einen
%   Rueckverweis dorthin bekommen und dafuer je einen Satz verlieren.
% =========================================================================


% ---------------------------- AB HIER KOPIEREN ---------------------------

\subsection{Analyserahmen und Baselines}
\label{sec:rahmen-baselines}

Von der Gesamtvarianz der Einsatzhäufigkeit liegen 92,5\,\% zwischen den
Stadtteilen (Abschnitt~\ref{sec:eda}). Ein Modell, das einen Stadtteil bereits
aus dem Training kennt, muss dessen Niveau deshalb nicht mehr erklären. Die
Aufteilung in Trainings- und Testmengen wird daher nicht entlang der Zeit,
sondern entlang der Stadtteile gezogen; Ziel der Validierung ist die
Generalisierung auf unbekannte Stadtteile. Sie entsteht in der Aufbereitung und
wird als Spalten \texttt{fold} und \texttt{ist\_holdout} in beide
Analysedateien geschrieben, ist also eine Eigenschaft des Datensatzes und nicht
der Algorithmen. Das Vorgehensmodell führt die Aufteilung folgerichtig unter
der Auswahl der Daten und nicht unter der
Modellierung.\footcite[23-26]{Chapman2000} Am Ende stehen fünf Folds über 29
Entwicklungsstadtteile, sechs zurückgehaltene Stadtteile und eine zweistufige
Messlatte, an der sich jedes spätere Verfahren messen lassen muss.

Geprüft wird, indem ganze Stadtteile zurückgehalten werden. Die 29
Entwicklungsstadtteile verteilen sich auf fünf Folds mit sechs, sechs, sechs,
sechs und fünf Stadtteilen; jeder ist genau einmal Testfall, und zwar mit allen
132 Monaten. Je Fold stehen damit 3.036 bis 3.168 Trainingszeilen und 660 bis 792
Testzeilen nebeneinander (Abbildung~\ref{abb:foldstruktur});
Quelltext~\ref{lst:foldmasken} zeigt, dass beide Mengen aus den Spalten der
Datei gebildet werden und kein Stadtteil in beiden steht. Ein Zeitschnitt
prüfte die Forschungsfrage nicht: Dort steht jeder Stadtteil in Training und
Test, und das Modell kennt sein Niveau bereits --- der Stadtteilmittelwert
allein erreicht ein Bestimmtheitsmaß von 0,888 (Abschnitt~\ref{sec:eda}). Für
genau diesen Fall empfehlen Roberts et al. eine räumlich geblockte Aufteilung,
wenn auf neue Orte vorhergesagt werden soll.\footcite[925]{Roberts2017} Was der
Rahmen dadurch nicht mehr misst, ist die Fortschreibung eines bereits bekannten
Stadtteils in die Zukunft.

Damit weicht die Untersuchung von der angemeldeten Fassung ab. Die zweite
Unterfrage nannte eine zeitreihengerechte Kreuzvalidierung, und ein früher
Entwurf sah ein Testfenster der letzten zwölf Monate vor. Der Wechsel auf den
stadtteilweisen Rahmen ist am 28.~Juli 2026 entschieden worden, nachdem die
Varianzzerlegung des Bestands vorlag. Er ist damit die Folge eines Befundes
über die Daten und keine Anpassung an ein Ergebnis; die Abwägung beider Rahmen
steht in Kapitel~\ref{sec:diskussion}.

Welcher Stadtteil in welchen Fold fällt, ist nicht zufällig. Die Zuteilung
folgt einer Ordnung, die zuerst nach der Zahl brand-dominierter Monate und bei
Gleichstand nach der Wohnbevölkerung sortiert; die so geordneten Stadtteile
werden reihum auf sechs Gruppen verteilt, von denen die erste das Hold-out
bildet (Quelltext~\ref{lst:folds}). Beide Kriterien beantworten je ein Problem.
Von 70 brand-dominierten Monaten liegen 35 in einem einzigen Stadtteil, und
ohne Stratifizierung enthielt in drei von vier Aufteilungen ein Fold keinen
einzigen Brand-Testfall; das Macro-F1 mittelte dann über eine Klasse, die im
Testfold nicht vorkam. Jetzt verteilen sich die Brandfälle mit 13, 9, 6, 3 und
2 Monaten auf die fünf Folds. Das zweite Kriterium verhindert, dass ein Fold
nur aus großen Stadtteilen besteht und die Streuung zwischen den Folds ein
Größeneffekt wäre. Ein Leckageproblem entsteht dabei nicht: Festgelegt wird
allein, welche Stadtteile gemeinsam getestet werden, nicht, was ein Modell über
sie lernt.

Sechs Stadtteile bleiben während der gesamten Entwicklung unberührt und werden
erst für die Schlussbewertung geöffnet. Sie sind nicht zufällig gezogen,
sondern die erste Gruppe derselben Ordnung, und deshalb kein verkleinertes
Abbild der Stadt: Auf sie entfallen 37 der 70 brand-dominierten Monate, ihr
Brandanteil liegt bei 4,67\,\% gegenüber 0,86\,\% in den
Entwicklungsstadtteilen. Der Stadtteil mit den meisten brand-dominierten
Monaten fällt nach dieser Regel zwangsläufig ins Hold-out. Ergebnisse auf dem
Hold-out sind daher nicht als Durchschnittsleistung über die Stadt zu lesen,
sondern als Probe unter erschwerten Bedingungen.

Eine einzelne Aufteilung von 29 Stadtteilen auf fünf Folds ist eine
willkürliche Gruppierung. Alle Referenzwerte entstehen deshalb aus 50 Läufen:
fünf Folds in zehn Wiederholungen, die innerhalb der Rangblöcke der
Zuteilungsordnung neu mischen und dabei Foldgrößen und Stratifizierung
erhalten; die erste Wiederholung reproduziert die Aufteilung der Datei
bitgenau. Berichtet wird die Streuung dieser zehn Wiederholungsmittel und nicht
die Streuung über alle 50 Läufe. Letztere fiele kleiner aus, obwohl ihr
dieselben 29 Stadtteile in zehn Anordnungen zugrunde liegen, und wäre damit zu
optimistisch.

Die Zahl der Zeilen überzeichnet zugleich, wie viel unabhängige Information
vorliegt. Die 132 Monate eines Stadtteils verhalten sich weitgehend wie
Wiederholungen derselben Beobachtung; der Intraklassenkorrelationskoeffizient
der Einsatzhäufigkeit beträgt 0,926 (Abschnitt~\ref{sec:eda}). Aus ihm und der
Zahl der Monate je Stadtteil ergibt sich ein Designeffekt von 122 und damit
eine effektive Stichprobe von rund 38 statt 4.620
Beobachtungen.\footcite[206]{Killip2004} Die Rechnung setzt gleich große
Gruppen voraus; diese Bedingung ist hier genau erfüllt, weil das Raster aus
Abschnitt~\ref{sec:merkmale} lückenlos 35 mal 132 Zeilen führt. Die Aufteilung
nach Stadtteilen trägt dem Rechnung, indem sie den Stadtteil und nicht die
Zeile als Einheit behandelt. Für die Auswertung folgt daraus, dass Unterschiede
zwischen Verfahren an der Streuung über die Wiederholungen zu beurteilen sind
und nicht an der Zahl der Zeilen.

Dass alle Verfahren dieselben Merkmale und dieselbe Aufteilung sehen, ist keine
Zusicherung, sondern eine Eigenschaft des Aufbaus. Die Fold-Zuteilung entsteht
in einem einzigen Aufruf und wird von dort in beide Analysedateien geschrieben;
Mengen- und Strukturstrang können konstruktiv keine unterschiedlichen
Aufteilungen sehen. Die Merkmalsliste steht genau einmal, in
\texttt{prep/config.py}, und wird von jedem Modellskript von dort gelesen
(Quelltext~\ref{lst:merkmalsliste}); eine zweite Stelle, an der ein
Merkmalssatz definiert wäre, gibt es nicht. Eine automatisierte Prüfung baut
die Aufteilung aus den Dateien nach und vergleicht sie zeilenweise mit den
gespeicherten Spalten. Ohne sie liefen bei einer veralteten Aufteilung alle
Verfahren auf falschen, aber untereinander identischen Mengen: Der Vergleich
bliebe fair und das Ergebnis wäre trotzdem falsch.

Beide Dateien werden abschließend auf modelltaugliche Datentypen gebracht ---
alle Merkmale als \texttt{float64}, Schlüssel und Zählgrößen als
\texttt{int64}, keine fehlenden Werte (Quelltext~\ref{lst:dtypes}). Der Aufwand
hat einen konkreten Anlass: Die Aggregation der Zensusdaten liefert
Medianeinkommen und Medianmiete als nullbare Ganzzahlen, und eine einzige
solche Spalte genügt, damit die Merkmalsmatrix beim Übergang in ein
NumPy-Feld den Typ \texttt{object} annimmt. Scikit-learn wandelt das still um,
XGBoost lehnt es ab. Der Fehler träte also erst beim dritten der drei zu
vergleichenden Verfahren auf und säße dann in der Aufbereitung und nicht im
Modell.

Die Messlatte ist zweistufig angelegt (Abschnitt~\ref{sec:zielgroessen}).
Stufe~1 kommt ohne ein einziges Merkmal aus: der Gesamtmittelwert der
Trainingsstadtteile für die Mengenzielgrößen, die häufigste Klasse für die
Einsatzart. Sie beantwortet, ob in den Merkmalen überhaupt Information steckt.
Stufe~2 ist die einfachste Form, die zur Datenform passt: ein verallgemeinertes
lineares Modell mit kanonischem Link, unpenalisiert angepasst. Beide Stufen
laufen über dieselben 50 Läufe, dieselbe Aufteilung und dieselben zwölf
Merkmale wie die späteren Vergleichsverfahren; die Baseline ist damit
Mitbewerber unter identischem Protokoll und nicht bloß ein Erwartungswert.
Maßgeblich ist Stufe~2 --- ein Verfahren, das nur Stufe~1 schlägt, hat nichts
gezeigt.

Für die Mengenzielgrößen ist Stufe~2 ein Poisson-Modell mit Offset
(Tabelle~\ref{tab:baselines-menge}, Quelltext~\ref{lst:poisson}). Die
logarithmierte Wohnbevölkerung geht mit fest auf eins gesetztem Koeffizienten
ein; geschätzt wird damit die Rate, hochgerechnet wird erst in der Vorhersage.
Drei Eigenschaften machen diese Wahl zur Messlatte und nicht zum Strohmann.
Über den Logarithmus als Link wirkt das Modell auf der Skala der Zielgröße
multiplikativ und bildet damit überproportionale Effekte ab; es erfasst
Krümmung, aber keine Wechselwirkungen --- und genau die finden Random Forest
und XGBoost konstruktionsbedingt. Es hat keinen frei wählbaren Parameter, weil
die Anpassung ohne Strafterm erfolgt; es gibt nichts zu optimieren und damit
auch nichts zu begründen. Und die zweite Mengenzielgröße entsteht aus derselben
Anpassung, geteilt durch die Wohnbevölkerung: Ein eigenes Ratenmodell wäre eine
zweite Spezifikation und gegenüber den Vergleichsverfahren nicht mehr dasselbe
Protokoll.

Für die Einsatzart ist Stufe~2 ein multinomiales Logit, ebenfalls ohne
Strafterm (Tabelle~\ref{tab:baselines-struktur},
Quelltext~\ref{lst:logit}). Es ist das Gegenstück zum Poisson-Modell: dieselbe
Modellklasse, derselbe kanonische Link, derselbe Verzicht auf einen freien
Parameter. Die vier Klassen gehen mit ausgeglichenen Gewichten ein; die
Ungleichverteilung bleibt damit eine Eigenschaft der Daten und wird nicht durch
Vervielfachen oder Entfernen von Zeilen verändert. Der Datensatz, auf dem
verglichen wird, ist in beiden Strängen derselbe.

Die Zähldaten sind deutlich überdispers; der Dispersionsindex liegt bei 62,8
(Abschnitt~\ref{sec:eda}). Die naheliegende Erweiterung wäre eine negative
Binomialverteilung, und sie war ursprünglich vorgesehen. Sie löst jedoch ein
Problem, das diese Baseline nicht hat: Überdispersion beschädigt beim
Poisson-Schätzer die Standardfehler, nicht die Konsistenz des bedingten
Mittelwerts, solange dieser richtig spezifiziert
ist.\footcite[5]{Gourieroux1981} Eine Referenz, die ausschließlich
Punktvorhersagen liefert, verwendet keine Standardfehler. Zugleich brächte die
negative Binomialverteilung mit dem Dispersionsparameter eine zusätzlich zu
schätzende Größe mit und wäre damit nicht mehr die einfachste Form, die zur
Datenform passt. Gemessen ist das Poisson-Modell außerdem die härtere
Messlatte: 33,98 gegenüber 37,27 RMSE.

Die Reduktion auf die einfacheren Varianten ist methodisch sauber, sie
vermeidet willkürliche Parameter und liefert zudem stärkere Vergleichswerte.
Tragend ist das zweite Argument: Ein Strafterm auf dem Vorgabewert der
verwendeten Software wäre ein willkürlich gesetzter Parameter, ein optimierter
eine zusätzlich zu begründende Wahl; die unpenalisierte Anpassung hat nichts zu
wählen. Das dritte Argument gilt allerdings nicht durchgehend. Im Mengenstrang
steigt die Messlatte von 37,27 auf 33,98 RMSE, im Strukturstrang sinkt sie von
0,314 auf 0,297 Macro-F1. Dass dieselbe Änderung den einen Strang härter und
den anderen weicher macht, ist zugleich das Argument dagegen, sie nachträglich
nach dem Ergebnis auszuwählen.

Drei Eigenschaften dieser Referenzwerte werden leicht als Fehler gelesen.
Negative Bestimmtheitsmaße sind hier richtig und aussagekräftig: Wer für einen
unbekannten Stadtteil den Gesamtdurchschnitt vorhersagt, liegt schlechter als
dessen eigener Mittelwert --- genau diese Lücke sollen die Strukturmerkmale
schließen. Auf der Rate ist das Bestimmtheitsmaß darüber hinaus kein geeignetes
Hauptmaß, weil es gegen den Mittelwert der jeweiligen Testdaten misst und die
Rate zwischen den Stadtteilen um den Faktor 32 streut; in einem Fold fällt es
auf $-0{,}920$, obwohl das Poisson-Modell die Marke der Stufe~1 beim RMSE in
jedem Fold schlägt. Für die Rate sind deshalb RMSE und mittlerer absoluter
Fehler maßgeblich und das Bestimmtheitsmaß nur nachrichtlich
(Abschnitt~\ref{sec:modellbewertung}). In der Klassifikation schließlich liegt
die Trefferquote der Stufe~2 unter der von Stufe~1, während das Macro-F1
steigt; maßgeblich ist Macro-F1, weil die Trefferquote von der häufigsten
Klasse bestimmt wird.

Eine naheliegende Referenz fehlt bewusst. Eine Fortschreibung des
Vormonatswerts ließe sich in diesem Rahmen nicht bilden, weil sie die eigene
Vergangenheit des Teststadtteils verwendete --- und ein zurückgehaltener
Stadtteil hat im Training keine. Aus demselben Grund gehen die drei
Verlaufsmerkmale aus Abschnitt~\ref{sec:merkmale} nicht in die Modelle ein.
Gesamtmittelwert und saisonaler Durchschnitt kämen als merkmalsfreie Referenzen
beide in Frage, fallen nachgerechnet aber praktisch zusammen und sind damit
faktisch eine triviale Referenz und nicht zwei. Die Anforderung an eine
Referenz, die nicht linear in der Zielgröße ist, erfüllt deshalb erst Stufe~2.

Tabelle~\ref{tab:steckbrief} fasst zusammen, was die Modellierung
übernimmt. Die letzte Zeile ist zugleich der Grund, warum die Angaben dieses
Kapitels nachprüfbar sind: 19 automatisierte Prüfungen laufen gegen die
erzeugten Dateien und nicht gegen den Code, der sie erzeugt hat. Sie prüfen
unter anderem, dass das Panel rechteckig und der Zeitraum festgesetzt ist, dass
kein Stadtteil zugleich Trainings- und Testfall ist
(Quelltext~\ref{lst:test-folds}), dass jeder Fold den vollen Zeitraum abdeckt,
dass die Verlaufsmerkmale keinen Gegenwartsbezug haben und dass keine der zehn
gesperrten Ergebnisvariablen in einer der beiden Dateien steht.

% ---------------------------- BIS HIER KOPIEREN --------------------------
